% =====================================================================
\section{ШИМ: яркость светодиода и частота}
\label{les:pwm}
% =====================================================================
\lessonintro
  {понять, почему быстрое мигание кажется ровным светом, освоить ШИМ и плавно менять яркость светодиода; научиться писать свои функции}
  {урок~\ref{les:signals} (ШИМ, двоичные числа), урок~\ref{les:led} (светодиод, период и частота)}
  {«дышащий» светодиод, плавная волна на четырёх светодиодах, мелодия на пищалке}
  {схема из урока~\ref{les:led}; пассивная пьезопищалка (по желанию)}

\subsection{Можно ли светить вполсилы?}

Команда \code{digitalWrite()} умеет только две вещи: включить вывод (3,3~В) или выключить (0~В).
Светодиод либо горит в полную силу, либо не горит совсем. Как же сделать ночник, который светит
тускло? У ESP32-S3 нет выхода, который выдавал бы, скажем, 1,5~В. Зато есть глаз, который
легко обмануть. Проведём опыт.

\subsection{Опыт: мигаем всё быстрее}

Возьми схему со светодиодом на \pin{4} из прошлого урока и загрузи программу, которая каждые
6 вспышек сокращает время вдвое:

\codecap{Опыт с частотой мигания}
\codeinput{l06_flicker}

Наблюдай за светодиодом и за сообщениями в Serial Monitor:

\begin{table}[H]
\centering\small
\begin{tabular}{@{}cccl@{}}
\toprule
\code{delay()}, мс & Период $T$, мс & Частота $f$, Гц & Что видно\\
\midrule
500 & 1000 & 1 & спокойное мигание\\
125 & 250 & 4 & быстрое мигание\\
31 & 62 & 16 & дрожащий, мерцающий свет\\
15 & 30 & 33 & лёгкое мерцание, особенно краем глаза\\
7 & 14 & 71 & \textbf{ровный свет}, но слабее обычного\\
3 & 6 & 167 & ровный свет\\
\bottomrule
\end{tabular}
\caption{Что видит глаз при разной частоте мигания}
\end{table}

Удивительно: светодиод продолжает включаться и выключаться, но мы видим ровный свет --- примерно
вполовину яркости! Почему так происходит?

\subsection{Инерция зрения}

Глаз --- не фотоаппарат с мгновенным затвором. Светочувствительные клетки сетчатки реагируют на свет
не сразу и «остывают» тоже не сразу: каждое впечатление держится примерно $1/20$--$1/50$ секунды.
Это называется \textbf{инерцией зрения}. Если вспышки следуют чаще, они сливаются, и мозг видит
\emph{среднюю} яркость.

\begin{figure}[H]
\centering
\begin{tikzpicture}
\begin{groupplot}[
  group style={group size=1 by 3, vertical sep=12mm},
  width=0.9\textwidth, height=3.3cm,
  /pgf/declare function={
    tauE=0.012;
    ymx(\P)=1/(1+exp(-(\P/2)/tauE));
    md(\u,\P)=\u-\P*floor(\u/\P);
    sq(\u,\P)=(md(\u,\P)<\P/2)?1:0;
    eye(\u,\P)=(md(\u,\P)<\P/2) ? 1-ymx(\P)*exp(-md(\u,\P)/tauE)
                               : ymx(\P)*exp(-(md(\u,\P)-\P/2)/tauE);
  },
  ymin=-0.15, ymax=1.3, ytick={0,1}, yticklabels={0,1},
  xlabel={$t$, с}, xlabel style={font=\scriptsize,at={(1,0)},anchor=north west},
  tick label style={font=\scriptsize}, title style={font=\sffamily\small,yshift=-1mm},
  scaled x ticks=false, xticklabel style={/pgf/number format/fixed,/pgf/number format/use comma},
]
  \nextgroupplot[title={2 Гц: глаз успевает за каждой вспышкой}, xmin=0, xmax=1, samples=600, domain=0:1]
    \addplot[thin,esp] {sq(x,0.5)};
    \addplot[very thick,accent] {eye(x,0.5)};
  \nextgroupplot[title={16 Гц: глаз едва успевает --- видно мерцание}, xmin=0, xmax=0.5, samples=800, domain=0:0.5]
    \addplot[thin,esp] {sq(x,0.0625)};
    \addplot[very thick,accent] {eye(x,0.0625)};
  \nextgroupplot[title={100 Гц: вспышки сливаются в ровный свет половинной яркости}, xmin=0, xmax=0.1, samples=900, domain=0:0.1,
                 xtick={0,0.02,...,0.1}]
    \addplot[thin,esp] {sq(x,0.01)};
    \addplot[very thick,accent] {eye(x,0.01)};
    \addplot[dashed,gray] coordinates {(0,0.5) (0.1,0.5)};
\end{groupplot}
\end{tikzpicture}
\caption{Один и тот же прямоугольный сигнал (\textcolor{esp}{красная линия}) с разной частотой.
\textcolor{accent}{Синяя линия} --- упрощённая модель того, что «видит» глаз}
\end{figure}

\begin{fact}
Инерция зрения --- основа кино и мультфильмов. Кинокадры сменяются 24 раза в секунду, а каждый кадр
проецируется два-три раза, чтобы вспышки шли с частотой 48--72~Гц. Экран телефона обновляется
60--120 раз в секунду. А обычные лампы в розетке мерцают 100 раз в секунду --- это можно увидеть,
если посмотреть на них через камеру телефона в режиме замедленной съёмки.
\end{fact}

\subsection{ШИМ --- широтно-импульсная модуляция}

Теперь понятна идея: будем мигать \emph{так быстро}, что глаз не видит мигания, а яркость менять,
изменяя \emph{долю времени}, пока светодиод горит. Это и есть \textbf{ШИМ} --- широтно-импульсная
модуляция (по-английски PWM, \emph{pulse-width modulation}): мы меняем («модулируем») ширину импульсов.

Доля времени во включённом состоянии называется \textbf{коэффициентом заполнения} $D$ (англ. \emph{duty cycle}):
\[
  D = \frac{t_\text{вкл}}{T}\cdot 100\,\%, \qquad U_\text{среднее} = D \cdot 3{,}3~\text{В}.
\]

\begin{figure}[H]
\centering
\begin{tikzpicture}[x=0.9cm,y=0.75cm]
  \foreach \d/\name [count=\row] in {0.1/10\%,0.5/50\%,0.9/90\%}{
    \begin{scope}[yshift=-\row*1.9cm]
      \draw[->] (0,0) -- (13,0) node[right,font=\scriptsize]{$t$};
      \draw (0,0) -- (0,1.6);
      \node[left,font=\sffamily\small,align=right] at (-0.1,0.7) {$D=\name$};
      \foreach \k in {0,...,3}{
        \pgfmathsetmacro{\s}{\k*3}
        \pgfmathsetmacro{\e}{\k*3+3*\d}
        \draw[very thick,esp] (\s,0) -- (\s,1.3) -- (\e,1.3) -- (\e,0) -- (\s+3,0);
      }
      \draw[dashed,accent,thick] (0,1.3*\d) -- (12,1.3*\d) node[right=3mm,font=\scriptsize,text=accent]{средний уровень};
      % лампочка-индикатор яркости
      \shade[ball color=red!\fpeval{round(\d*100)}!black!30] (16.5,0.7) circle (0.45);
    \end{scope}
  }
  \draw[decorate,decoration={brace,amplitude=3pt}] (0,-1.9cm+1.1cm) -- (3,-1.9cm+1.1cm) node[midway,above=3pt,font=\scriptsize]{период $T$};
  \node[font=\sffamily\scriptsize] at (16.5,-0.3) {яркость};
\end{tikzpicture}
\caption{ШИМ с разным коэффициентом заполнения: чем шире импульсы, тем ярче светодиод}
\end{figure}

\subsection{ШИМ в ESP32-S3: модуль LEDC}

В ESP32-S3 ШИМ создаёт специальный блок \textbf{LEDC} (от англ. \emph{LED control} --- управление
светодиодами). Он переключает вывод сам, аппаратно, и процессору не нужно этим заниматься.
Нужно задать всего три вещи:
\begin{itemize}
  \item \textbf{вывод};
  \item \textbf{частоту} ШИМ --- сколько раз в секунду вывод включается. Для светодиодов берут 5000~Гц:
        это в 50 раз быстрее, чем может заметить глаз;
  \item \textbf{разрешение} в битах --- на сколько ступенек делится период. 8 бит дают $2^8 = 256$ ступенек:
        значения от 0 (всегда выключен) до 255 (всегда включён). Помнишь таблицу двоичных чисел из
        урока~\ref{les:signals}?
\end{itemize}

\begin{figure}[H]
\centering
\begin{tikzpicture}[x=0.55mm,y=1cm]
  \draw[->] (0,0) -- (275,0) node[right,font=\scriptsize]{ступеньки};
  \draw[very thick,esp] (0,0) -- (0,1) -- (64,1) -- (64,0) -- (256,0) -- (256,1);
  \foreach \x in {0,64,128,192,256} \draw (\x,0.05) -- (\x,-0.05) node[below,font=\scriptsize]{\x};
  \draw[decorate,decoration={brace,amplitude=4pt}] (0,1.15) -- (64,1.15) node[midway,above=4pt,font=\scriptsize]{\code{ledcWrite(pin, 64)}: 64 из 256 = 25\,\%};
  \draw[{Stealth[length=2mm]}-{Stealth[length=2mm]}] (0,-0.7) -- node[below,font=\scriptsize]{один период = 1/5000~с = 0,2~мс = 200~мкс} (256,-0.7);
\end{tikzpicture}
\caption{8-битный ШИМ: период разделён на 256 ступенек. Значение \code{duty} говорит, сколько из них вывод включён}
\end{figure}

В Arduino-ESP32 для этого есть две функции: \code{ledcAttach(pin, частота, биты)} --- вместо
\code{pinMode()}, и \code{ledcWrite(pin, duty)} --- вместо \code{digitalWrite()}.

\subsection{Светодиод с линейно меняющейся яркостью}

Сделаем «дыхание»: пусть значение \code{duty} равномерно растёт от 0 до 255, а потом так же
равномерно падает.

\codecap{«Дыхание» светодиода: \code{duty} меняется линейно}
\codeinput{l06_breathe}

\begin{figure}[H]
\centering
\begin{tikzpicture}
\begin{axis}[
  width=0.8\textwidth, height=4.4cm,
  xlabel={время, с}, ylabel={duty},
  xmin=0, xmax=4.2, ymin=0, ymax=280, ytick={0,128,255},
  grid=major, grid style={gray!25},
  tick label style={font=\scriptsize}, label style={font=\small},
]
  \addplot[very thick,esp] coordinates {(0,0) (1.02,255) (2.04,0) (3.06,255) (4.08,0)};
\end{axis}
\end{tikzpicture}
\caption{Как меняется \code{duty} в программе: 256 шагов по 4~мс вверх, столько же вниз --- «пила» с периодом около 2~с}
\end{figure}

Здесь два цикла \code{for}: первый считает вверх (\code{d++}), второй --- вниз (\code{d--}), начиная с 255
и пока \code{d >= 0}. Каждые 4~мс \code{duty} меняется на единицу. За это время ШИМ успевает сделать
$5000 \cdot 0{,}004 = 20$ периодов --- глаз видит не мигание, а плавное изменение яркости.

\subsection{Почему «линейно» выглядит нелинейно}

Запусти пример и присмотрись. Светодиод очень быстро становится ярким, а потом долго почти не
меняется --- хотя \code{duty} растёт равномерно! Дело снова в глазе.

Глаз замечает не саму яркость, а \emph{во сколько раз} она изменилась. Разница между 1 и 2 лампочками
огромна (вдвое!), а между 100 и 101 --- незаметна. Поэтому в тёмной части нужны очень маленькие шаги,
а в яркой --- большие. Исправляется это \textbf{гамма-коррекцией}: вместо равномерного роста \code{duty}
возводим «желаемую яркость» в степень 2,2:
\[
  \text{duty} = 255 \cdot \left(\frac{\text{яркость в \%}}{100}\right)^{2{,}2}.
\]

\begin{figure}[H]
\centering
\begin{tikzpicture}
\begin{groupplot}[
  group style={group size=2 by 1, horizontal sep=18mm},
  width=0.47\textwidth, height=5cm,
  grid=major, grid style={gray!25},
  tick label style={font=\scriptsize}, label style={font=\small},
  title style={font=\sffamily\small},
  legend style={font=\scriptsize,at={(0.03,0.97)},anchor=north west},
]
  \nextgroupplot[title={что подаём на вывод},
    xlabel={желаемая яркость, \%}, ylabel={duty (0..255)}, xmin=0, xmax=100, ymin=0, ymax=260]
    \addplot[very thick,gray,dashed,domain=0:100] {2.55*x};
    \addlegendentry{линейно}
    \addplot[very thick,esp,domain=0:100,samples=60] {255*(x/100)^2.2};
    \addlegendentry{с гаммой 2,2}
  \nextgroupplot[title={что видит глаз},
    xlabel={желаемая яркость, \%}, ylabel={кажущаяся яркость, \%}, xmin=0, xmax=100, ymin=0, ymax=105]
    \addplot[very thick,gray,dashed,domain=0:100,samples=80] {100*(x/100)^(1/2.2)};
    \addlegendentry{линейно: «рывок» в начале}
    \addplot[very thick,esp,domain=0:100] {x};
    \addlegendentry{с гаммой: ровно}
\end{groupplot}
\end{tikzpicture}
\caption{Гамма-коррекция: на вывод подаём «изогнутую» кривую, чтобы глаз увидел ровное изменение яркости}
\end{figure}

\subsection{Учимся программировать: свои функции}

Формулу гамма-коррекции придётся использовать много раз: для каждого значения яркости. Чтобы не
переписывать её снова и снова, оформим её как \textbf{свою функцию}. Функция --- это «машинка»: у неё есть имя,
на вход ей передают данные (\textbf{параметры}), а на выходе она \textbf{возвращает} результат. Написал один
раз --- пользуешься сколько угодно, как новым словом в языке.

\begin{figure}[H]
\centering
\begin{tikzpicture}
  \node[draw,thick,fill=blkcpu,rounded corners=3mm,minimum width=46mm,minimum height=18mm,align=center,font=\sffamily] (f) {функция \code{gammaDuty}\\[1mm] \scriptsize\code{255 * pow(percent / 100.0, 2.2)}};
  \draw[arr,very thick] ($(f.west)+(-28mm,0)$) node[left,font=\ttfamily]{50} -- node[above,font=\sffamily\scriptsize]{параметр \code{percent}} (f.west);
  \draw[arr,very thick] (f.east) -- node[above,font=\sffamily\scriptsize]{\code{return}} ($(f.east)+(24mm,0)$) node[right,font=\ttfamily]{56};
\end{tikzpicture}

\vspace{5mm}
\begin{tikzpicture}[tok/.style={font=\ttfamily\large,inner sep=1pt},ann/.style={font=\sffamily\scriptsize,align=center,text=black}]
  \node[tok,text=esp] (a) at (0,0) {uint32\_t};
  \node[tok,text=accent,base right=0.5em of a] (b) {gammaDuty};
  \node[tok,base right=0pt of b] (p) {(};
  \node[tok,text=okgreen!60!black,base right=0pt of p] (c) {int percent};
  \node[tok,base right=0pt of c] (e) {) \{ \ldots\ \}};
  \draw[esp,thick] (a.south) -- ++(0,-0.35) node[ann,below] {тип результата\\ (\code{void} --- без результата)};
  \draw[accent,thick] (b.south) -- ++(0,-1.1) node[ann,below] {имя};
  \draw[okgreen!60!black,thick] (c.south) -- ++(0,-0.35) node[ann,below] {параметр:\\ тип и имя};
\end{tikzpicture}
\caption{Функция получает яркость в процентах и возвращает значение \code{duty}: 50\,\% → 56}
\end{figure}

\begin{itemize}
  \item Слово \code{return} («вернуть») отдаёт результат туда, откуда функцию вызвали. Поэтому
        \code{ledcWrite(LED\_PIN, gammaDuty(p))} сначала вычислит \code{gammaDuty(p)}, а потом подставит результат.
  \item Функции \code{setup()} и \code{loop()} имеют тип \code{void} («пустота»): они ничего не возвращают.
  \item Почему в формуле \code{100.0}, а не \code{100}? Помнишь целочисленное деление из урока~\ref{les:ide}?
        \code{50 / 100} даст 0, а \code{50 / 100.0} --- правильные 0,5.
\end{itemize}

\begin{tip}
Хорошая функция делает \emph{одно} дело, и её имя говорит, какое: \code{gammaDuty()}, \code{blinkTwice()},
\code{isButtonPressed()}. Тогда программа читается почти как обычный текст.
\end{tip}

\codecap{«Дыхание» с гамма-коррекцией}
\codeinput{l06_breathe_gamma}

Функция \code{pow(x, y)} возводит $x$ в степень $y$. Сравни два примера: теперь светодиод
разгорается и гаснет плавно по всей длине.

\begin{fact}
Глаз способен видеть и при свете луны, и на солнечном пляже, где света в миллион раз больше. Такой
огромный диапазон возможен как раз потому, что зрение реагирует на \emph{отношения}, а не на разницу.
Слух устроен так же, поэтому громкость звука измеряют в децибелах --- логарифмической шкале.
Этот закон восприятия называют законом Вебера--Фехнера.
\end{fact}

\subsection{Какую частоту ШИМ выбрать}

Частота и разрешение ШИМ связаны: блок LEDC работает от внутреннего «метронома» 80~МГц, и за один
период ШИМ он должен успеть отсчитать все ступеньки. Чем выше частота, тем меньше ступенек помещается
в период.

\begin{table}[H]
\centering\small
\begin{tabular}{@{}lll@{}}
\toprule
Применение & Частота & Разрешение \\
\midrule
Светодиоды & \SI{5}{kHz} & 8--13 бит \\
Сервопривод (урок~\ref{les:adc}) & \SI{50}{Hz} & 14 бит \\
Моторчик через драйвер & \SIrange{10}{20}{kHz} & 10 бит \\
Звук (пищалка) & \SIrange{100}{5000}{Hz} & 8 бит \\
\bottomrule
\end{tabular}
\caption{Типичные настройки ШИМ}
\end{table}

\begin{tip}
Если выбрать для светодиода частоту всего 50--100~Гц, глаз не заметит мерцания, если смотреть прямо.
Но стоит быстро перевести взгляд или снять светодиод на камеру телефона --- появятся полосы и «пунктир».
Поэтому для света берут тысячи герц.
\end{tip}

\subsection{ШИМ и звук}

Звук --- это колебания воздуха. Если подключить к выводу \textbf{пьезопищалку} и включать-выключать
его, например, 440 раз в секунду, пластинка внутри пищалки будет колебаться с той же частотой, и мы
услышим ноту \emph{ля} первой октавы. Частота определяет \textbf{высоту звука}: чем больше герц, тем выше нота.

\begin{figure}[H]
\centering
\begin{tikzpicture}[x=1.4cm,y=0.7cm]
  \foreach \n/\f/\per/\yy in {{до (262 Гц)}/262/1.0/0,{до следующей октавы (523 Гц)}/523/0.5/-2.4}{
    \begin{scope}[yshift=\yy cm]
      \draw[->] (0,0) -- (8.3,0) node[right,font=\scriptsize]{$t$};
      \node[anchor=west,font=\sffamily\small] at (0,1.8) {\n};
      \pgfmathsetmacro{\N}{int(8/\per)-1}
      \foreach \k in {0,...,\N}{
        \draw[very thick,okgreen!70!black] (\k*\per,0) -- (\k*\per,1.2) -- (\k*\per+\per/2,1.2) -- (\k*\per+\per/2,0) -- (\k*\per+\per,0);
      }
    \end{scope}
  }
\end{tikzpicture}
\caption{Нота на октаву выше --- это ровно вдвое большая частота}
\end{figure}

Для звука используют функцию \code{ledcWriteTone(pin, частота)}: она сама настраивает ШИМ с
коэффициентом заполнения 50\,\% и нужной частотой. Подключи \emph{пассивную} пищалку (без встроенного
генератора) к \pin{15}:

\begin{twoview}{Пьезопищалка на выводе GPIO15}{fig:buzzer}
\begin{tikzpicture}[bb]
  \bbBoard{24}
  \bbDevKit{29}{25}
  \bbBuzzer{10}{12}
  \draw[wire=worange] (dk-15) -- (27.5,18) -- (27.5,23.5) -- (8,23.5) -- (8,14) -- (8,12) -- (10,12);
  \draw[wire=wblack] (12,14) -- (12,16);
  \draw[wire=wblack] (dk-GND) -- (26.5,4) -- (26.5,16) -- (24,16);
\end{tikzpicture}
\nextview
\begin{circuitikz}
  \draw (0,0) node[left]{\pin{15}} to[short,o-] (0.5,0) -- (1.5,0) to[buzzer,l=пищалка] (1.5,-2) node[ground]{};
  \node[font=\sffamily\scriptsize] at (1.3,-0.3) {+};
\end{circuitikz}
\end{twoview}

\codecap{Гамма «до-ре-ми» на пищалке}
\codeinput{l06_melody}

\begin{table}[H]
\centering\small
\begin{tabular}{@{}*{8}{c}@{}}
\toprule
до & ре & ми & фа & соль & ля & си & до\\
\midrule
262 & 294 & 330 & 349 & 392 & 440 & 494 & 523\\
\bottomrule
\end{tabular}
\caption{Частоты нот первой октавы, Гц. Каждая следующая октава --- вдвое выше}
\end{table}

\subsection{Практика: плавная волна}

Вернёмся к схеме из четырёх светодиодов (рис.~\ref{fig:led4}). Теперь у каждого будет свой ШИМ, и яркость
будет меняться по синусоиде --- той самой волне из урока~\ref{les:signals}. Если сдвинуть волну каждого
следующего светодиода на четверть периода, по ряду побежит мягкая световая волна.

\begin{figure}[H]
\centering
\begin{tikzpicture}
\begin{axis}[
  width=0.85\textwidth, height=4.6cm,
  xlabel={время, с}, ylabel={яркость, \%},
  xmin=0, xmax=2, ymin=0, ymax=105,
  grid=major, grid style={gray!25},
  tick label style={font=\scriptsize}, label style={font=\small},
  legend style={font=\scriptsize,at={(0.5,1.03)},anchor=south,legend columns=4,draw=none},
  trig format plots=rad, samples=120, domain=0:2,
]
  \addplot[very thick,red!80!black]    {50*(sin(2*pi*x)+1)};        \addlegendentry{GPIO4}
  \addplot[very thick,yellow!70!black] {50*(sin(2*pi*(x-0.25))+1)}; \addlegendentry{GPIO5}
  \addplot[very thick,green!60!black]  {50*(sin(2*pi*(x-0.5))+1)};  \addlegendentry{GPIO6}
  \addplot[very thick,blue!70]         {50*(sin(2*pi*(x-0.75))+1)}; \addlegendentry{GPIO7}
\end{axis}
\end{tikzpicture}
\caption{Яркость четырёх светодиодов: одна и та же волна, сдвинутая на четверть периода}
\end{figure}

\codecap{Световая волна}
\codeinput{l06_wave}

\begin{summary}
\begin{itemize}[leftmargin=*]
  \item Из-за инерции зрения вспышки чаще $\approx$50 раз в секунду сливаются в ровный свет.
  \item ШИМ меняет долю времени, пока вывод включён (коэффициент заполнения), и этим --- среднюю яркость.
  \item \code{ledcAttach(pin, freq, bits)} + \code{ledcWrite(pin, duty)}; при 8 битах \code{duty} от 0 до 255.
  \item Глаз воспринимает яркость нелинейно, поэтому для плавности нужна гамма-коррекция.
  \item Частота определяет, заметно ли мерцание, а для звука --- высоту ноты.
  \item Своя функция: тип результата, имя, параметры; \code{return} возвращает результат.
\end{itemize}
\end{summary}

\begin{quiz}
\begin{enumerate}
  \item Почему при частоте мигания 100~Гц мы видим ровный свет?
  \item Чему равен коэффициент заполнения, если вывод включён 150~мкс из периода 200~мкс?
  \item Какое значение \code{duty} соответствует 50~\% при 10-битном разрешении?
  \item Какая частота у ноты «ля» второй октавы, если у «ля» первой --- 440~Гц?
\end{enumerate}
\end{quiz}

\begin{tasks}
\begin{enumerate}
  \item Сделай «дыхание» в два раза медленнее, не меняя \code{delay()}. Подсказка: шаг.
  \item Напиши функцию \code{void blinkTimes(int n)}, которая мигает светодиодом \code{n} раз, и вызови её с разными числами.
  \item Попробуй функцию \code{ledcFade(pin, от, до, мс)}: она меняет яркость плавно и полностью сама.
  \item Сыграй на пищалке «В траве сидел кузнечик» или любую другую мелодию.
  \item Установи частоту ШИМ 30~Гц и посмотри на «дыхание». Что изменилось и почему?
\end{enumerate}
\end{tasks}
